\scanpage{269}
\mspage{10}
\margin{n° 227}
2. \uline{Systèmes linéaires}.

\uline{Addition au chapeau} : \uline{On suppose} $V$ \uline{complète}.

Soit $D$ un diviseur sur $V$. Les fonctions $z$ telles que $(z)\geq-D$ forment évidemment un espace vectoriel sur $k$, qui sera noté $L(D)$. Si $D_1\geq D_2$, on a $L(D_1)\supset L(D_2)$. La dimension de $L(D)$ sur $k$ sera notée $\ell(D)$. Le théorème qui suit montre que cette dimension est finie si $V$ est complète. Remarquons qu'en tout cas, $L(0)=k$, d'après § 1, prop.7.

THÉORÈME 1. -- \uline{Soit} $V$ \uline{complète}. \uline{Si} $D$ \uline{est un diviseur sur} $V$, \uline{alors} $L(D)$ \uline{est de dimension finie sur} $k$.

\uline{Dém}. La démonstration se fera par induction sur la dimension $r$ de $V$. Remarquons qu'il suffit de démontrer le théorème pour les diviseurs positifs. On peut aussi supposer $V$ normale, car on peut identifier les diviseurs sur $V$ et sur $V^*$.

Soit $V$ une courbe. Dans ce cas, l'absence de singularités de dimension $0$ implique que $V$ est normale. Les diviseurs premiers s'identifient aux points de $V$. Il suffit de montrer que si le théorème est vrai pour un diviseur $D$, il l'est aussi pour $D+Q$, où $Q$ est un point sur $V$. En fait, on démontrera que $\ell(D+Q)\leq\ell(D)+1$. (On part de $D=0$). Supposons qu'il existe deux fonctions $z$ et $w\in L(D+Q)$ qui sont linéairement indépendantes mod $L(D)$. Soit $n$ la multiplicité de $Q$ dans $D$. Puisque $z$ et $w$ ne sont pas dans $L(D)$, ces fonctions doivent avoir un pôle d'ordre $n+1$ sur $Q$, et par conséquent, $zw^{-1}$ est une unité à $Q$. Sous la place induite par $Q$ sur le corps de fonctions, $zw^{-1}$ prend une valeur constante $a\in k$. Alors $zw^{-1}-a$ possède un zéro sur $Q$ et $w(zw^{-1}-a)=z-aw$ un pôle d'ordre au plus

\mspage{268}
